\section{Introduction}
\label{sec:intro}

LLM users routinely encounter answers that are neither right nor an honest refusal: the model produces a plausible-sounding response that fails to engage with the question's actual demand. Standard honesty benchmarks score correctness versus abstention~\citep{lin2022truthfulqa, yin2023selfaware, cheng2024dontknow} and miss this third regime---a \emph{loophole} answer that nominally satisfies the surface request while violating its intent.

{\bf Why this matters.} The closest existing work, \citet{choi2025loopholes}, gives the first direct measurement of loophole \emph{exploitation} (not just judgment), but only for \emph{goal-conflict} settings where the model has the knowledge and its incentives conflict with compliance. Yet most everyday queries hit a different regime: the model simply does not know the answer. The honesty literature has framed this as the \emph{laziness $\leftrightarrow$ hallucination} continuum~\citep{zhao2024curriculum, simhi2025choke}, but neither extreme captures the loophole-shaped middle---the model says \emph{something} that engages the question without admitting that it does not know.

{\bf Gap.} {\bf Are knowledge-conflict loopholes real?} No paired-prompt benchmark exists for the case where the loophole comes from \emph{not knowing}. Existing knowledge-boundary datasets~\citep{yin2023selfaware, cheng2024dontknow, simhi2025choke} score refusal or correctness on unanswerable items, but they cannot \emph{causally} attribute behavior to knowability because they lack a matched knowable control. Without that control, every wrong answer is a hallucination by default, and the question of whether the model is \emph{strategically} dodging cannot even be asked.

{\bf Our approach.} We construct \kcl, a paired knowledge-conflict loophole benchmark of 20 templates, each instantiated as a topic-matched (knowable, unknowable) pair (e.g., the height of the Eiffel Tower vs. the height of the Yangshan Lighthouse). We run five LLMs spanning four capability tiers---frontier closed (\gptmini, \haiku), large open (\llamabig), small open (\llamasmall), and long-chain-of-thought (\rseventy)---collect 600 free-form generations and 600 multiple-choice intent probes, judge them with a four-category rubric, and define an outcome as \unsat if it is a confident wrong answer or a verbalized dodge. We complement \kcl with two external probes: a single-model replication of the \citet{choi2025loopholes} scalar implicature task (a methodological anchor to published numbers) and a \selfaware~\citep{yin2023selfaware} probe on genuinely unanswerable items to test whether the dodge mode emerges in a different question regime.

{\bf Quantitative preview.} Every model has a higher unsatisfactory rate on unknowable than knowable items, with per-pair deltas of $+0.13$ to $+0.28$. A binomial GLM with cluster-robust SEs on \texttt{template\_id} gives an odds ratio of $3.24$ ($95\%$ CI $[1.21, 8.68]$, $p = 0.020$) for knowability. The mechanism, however, is overwhelmingly \emph{confident hallucination} (\directans $\wedge$ \textsc{wrong} at $25$--$38\%$), not verbalized reinterpretation (\dodge at $\le 5\%$). On the \selfaware probe the mechanism inverts: $87$--$97\%$ dodge, $\le 10\%$ refusal. Long-chain-of-thought reasoning roughly halves the \kcl unsatisfactory rate (OR $0.50$, $p = 0.012$ vs. the small-open baseline), consistent with the goal-conflict finding of \citet{choi2025loopholes}. The single-model scalar replication on \gptmini reproduces the published rate ($0.583$ vs. $\sim\!0.62$) and the published direction of the \texttt{scalar\_max} drop ($0.312$).

{\bf Contributions.}
\begin{itemize}[leftmargin=*,itemsep=0pt,topsep=0pt]
    \item We introduce \kcl, the first paired-prompt benchmark for knowledge-conflict loophole behavior, and release a four-category judge rubric that operationalizes \emph{obviously unsatisfactory} as either confident-wrong or verbal dodge.
    \item We measure five LLMs on \kcl and show that knowability triples the odds of an unsatisfactory response (GLM OR $3.24$, $p = 0.020$), with two of five models reaching Bonferroni-corrected per-model significance.
    \item We complement \kcl with a \selfaware probe and a \citet{choi2025loopholes} scalar replication, and use the contrast to propose a two-mode loophole taxonomy: implicit hallucination loopholes on factual questions versus verbal dodge loopholes on underspecified ones.
    \item We document that long-chain-of-thought reasoning roughly halves the knowledge-conflict unsatisfactory rate ($\rseventy$, OR $0.50$, $p = 0.012$), extending the \citet{choi2025loopholes} finding from goal conflict to knowledge conflict.
\end{itemize}

The remainder of the paper is organized as follows. \Secref{sec:related} situates our work in the loophole, hallucination, and honesty literatures. \Secref{sec:methodology} describes the \kcl construction, judging pipeline, and statistical plan. \Secref{sec:results} presents the main effects, fine-grained breakdowns, intent-versus-behavior gap, and the \selfaware/\choi companion probes. \Secref{sec:discussion} develops the two-mode taxonomy and discusses limitations. \Secref{sec:conclusion} closes with future work.
